\documentclass[11pt]{article}

\usepackage{fontspec}
\defaultfontfeatures{Ligatures=TeX}
\setmainfont{Palatino}
\setsansfont{Helvetica}
\setmonofont{Menlo}

\usepackage{kotex}
\setmainhangulfont{NanumMyeongjo}

\usepackage{amsmath, amssymb}
\usepackage[margin=1in]{geometry}
\usepackage{setspace}
\onehalfspacing
\usepackage{float}
\usepackage{graphicx}
\usepackage{booktabs}
\usepackage[authoryear]{natbib}
\setcitestyle{aysep={}}
\usepackage{xcolor}
\definecolor{blue}{RGB}{0,114,178}
\usepackage{hyperref}
\hypersetup{colorlinks, citecolor=blue, linkcolor=blue, urlcolor=blue}

\begin{document}

\title{Two Routines, Not a Chain: Confirmation Hearing Opposition and the Allocation of Audit Scrutiny in the Korean National Assembly}
\author{KNA Research Agents (AI-generated) \\ Experimental Output - kna-research-agents.com}
\date{\today}
\maketitle
\thispagestyle{empty}

\begin{abstract}
\noindent Legislators who publicly oppose a ministerial nominee at a confirmation hearing might be expected to scrutinize that ministry more intensively afterward. Position-taking accounts of legislative oversight predict such carry-over; accounts of confirmation hearings as party theater predict none. I adjudicate these opposing predictions with a pre-registered difference-in-differences design that links 7.4 million committee question-and-answer dyads from the Korean National Assembly to hand-coded hearing rosters and party-based opposition status for three cohorts of ministerial nominees, with a speech-based coding of individually expressed opposition as a robustness check. Opposing a nominee does not raise the confirmed ministry's share of a legislator's questions at the following annual state audit. The estimate is statistically equivalent to zero at the threshold declared before estimation, placebo agencies move identically, an apparent baseline attention gap dissolves under committee fixed effects, and a within-opposition intensity slope, which the achievable sample cannot resolve, is bounded below half the pre-set bar. Confirmation conflict and audit allocation appear to be separate institutional routines rather than links in a causal chain.
\end{abstract}

\bigskip
\noindent\textbf{Keywords:} legislative oversight, confirmation hearings, parliamentary questions, equivalence testing, South Korea

\bigskip
\setcounter{page}{0}
\clearpage

\section{Introduction}

Legislatures oversee the executive through a sequence of distinct institutional venues \citep{mccubbins1984, aberbach1990}. They screen the personnel who will lead ministries, and they later scrutinize what those ministries do. An assumption connecting the two venues runs through both scholarship and institutional reform debate: the conflict generated when a nominee is screened should reverberate through the scrutiny that follows, because the legislators who fought the appointment acquired a stake in the ministry they failed to keep out of the nominee's hands. The assumption is consequential. In South Korea, recurring proposals to reorganize the National Assembly's confirmation hearing system, including proposals to divide the hearing into separate policy and ethics tracks, rest in part on beliefs about what the hearing contributes to the Assembly's subsequent oversight of the ministry \citep{choi2008, yoon2020}. Whether confirmation conflict in fact carries into later oversight behavior, however, has not been established in any national setting.

Two literatures make opposite predictions about this carry-over, and they make them about the same measurable quantity. The first extends the study of confrontation in legislative oversight. \citet{eldes2023} show that congressional oversight hearings mix information gathering with confrontation and that confrontation tracks partisan alignment with the administration. If confrontation at a confirmation hearing is individual position taking with reputational returns \citep{mayhew1974}, then a legislator who staked a public position against a nominee has invested in that ministry's failure. Position-taking continuity predicts that the investment pays dividends in attention: the confirmed ministry should command a larger share of that legislator's oversight questions after the hearing than before. The second literature is the Korean scholarship on confirmation hearings, which has consistently characterized hearing behavior as an expression of executive-legislative and partisan relations rather than of individual conviction \citep{choi2008, jeon2018, yoon2020}. On this account, opposing a nominee is a party-assigned role for the hearing day. The role is discharged when the hearing adjourns, and the legislator's oversight portfolio reverts to the jurisdiction of the standing committee. Continuity predicts a positive shift in attention toward the contested ministry; the party-theater account predicts none.

Despite mature literatures on confirmation politics, on parliamentary questioning, and on legislative oversight, there exists, to my knowledge, no study in any country that tests whether the positions legislators take at confirmation hearings carry into those same legislators' subsequent oversight behavior. The fragments sit in three separate bodies of work that have not been connected. Confirmation studies ask why nominees pass or fail \citep{yoon2020}. Oversight studies ask how legislators behave once a hearing is under way \citep{eldes2023}. The parliamentary questions literature asks how parties allocate scrutiny across issues \citep{martin2011, senninger2017}. None follows the individual legislator across the boundary between the appointment fight and the oversight venue that comes after it. The gap matters because the two theoretical traditions that could fill it disagree, and because a reform debate proceeds as though the answer were known.

This article closes the gap with data from the Korean National Assembly, a setting unusually well suited to the test. Ministerial nominees face mandatory confirmation hearings before the standing committee that holds jurisdiction over their ministry. Every autumn, the same standing committees conduct the state audit, a calendar-fixed burst of police-patrol oversight in the sense of \citet{mccubbins1984}, in which legislators question ministries and agencies directly; the audit falls five to ten months after the spring and early-summer hearing cohorts within the Assembly's rigid institutional time structure \citep{ka2025}. Because the same legislators question the same ministries in consecutive audits, the design can compare each legislator's allocation of audit attention before and after the hearing at which the legislator took a position on the nominee. I draw on a corpus of 9.9 million committee speech acts and 7.4 million question-and-answer dyads spanning the 16th through 22nd Assemblies, hand-code hearing rosters and party-based opposition status for three cohorts of ministerial nominees heard in 2020 to 2021, 2022, and 2023, and estimate a difference-in-differences model of the confirmed ministry's share of each legislator's audit questions; a second coding drawn from each legislator's own hearing speech probes the individually expressed construct. The primary threshold, a shift of five percentage points, and the falsification criteria were declared before estimation, following the logic of arguing for a negligible effect \citep{rainey2014}, and the placebo machinery follows the equivalence approach of \citet{hartman2018}.

The answer is a null at the margin the design was built to measure, and a bound at the one margin it cannot resolve. Legislators who opposed a nominee do not shift their audit questioning toward the confirmed ministry. The estimated change is small, statistically indistinguishable from zero, and formally equivalent to zero at the pre-registered threshold; the minimum detectable effect binds well below that threshold, so an effect of the size the continuity hypothesis predicts would have been detected with at least 80 percent probability. Placebo agencies, which no nominee headed, move identically, and the placebo estimate is equivalent to zero at even half the registered margin. An apparent baseline gap, in which opposed legislators already devoted more of their audit questions to the ministry before the hearing, dissolves under committee fixed effects: it reflects the composition of committee jurisdictions, driven almost entirely by one committee whose witness list is effectively a single ministry. Within opposition questioners, the intensity of hearing engagement is associated with a positive but small reallocation that the achievable sample cannot resolve; that margin is reported as inconclusive within a pre-set bound, with the confidence bound sitting below half the bar for a diluted continuity effect. A separate cohort of nominees whose hearings spanned a full change of government between audits shows the same null, suggesting the result is not an artifact of one partisan configuration.

The article makes four contributions. Empirically, it adjudicates a contradiction between two literatures and finds that neither positive mechanism operates at the individual level: audit allocation follows committee jurisdiction, and confirmation conflict leaves no measurable trace on it. This narrows the scope of position-taking continuity arguments and disciplines the party-theater reading, which had itself predicted a baseline attention pattern that the data attribute to composition. Methodologically, it offers a template for informative null results in legislative studies: pre-registered thresholds, binding minimum detectable effects, formal equivalence tests, and two self-corrections executed by pre-committed decision rules rather than post hoc judgment. For institutional reform, the Korean debate over restructuring confirmation hearings has assumed that hearing conflict has an oversight externality at the audit; the evidence here bounds that externality, equivalent to zero at the pre-registered threshold at the main margin and confined below half the pre-set bar at the intensity margin the design cannot fully resolve. On scope, finally, the article documents the limits of the claim with care: the attention literature predicts that confrontation-generated attention decays within weeks to months \citep{birkland1998, walgrave2007, vliegenthart2016}, so a short-range carry-over could exist and be fully discharged before October; the test here concerns the audit, the venue the carry-over claim itself names.

The next section derives the competing hypotheses and pre-registered thresholds from the oversight, confirmation, and political-attention literatures; the sections that follow describe the data and identification strategy, report the main estimates, equivalence tests, and secondary margins, and draw out the implications for theory and for the reform debate.

\section{Literature and Theory}

\subsection{Position Taking and the Continuity of Confrontation}

The study of legislative oversight long centered on institutional design, on the choice between centralized patrols and decentralized alarms and on the committee structures that carry scrutiny \citep{mccubbins1984, aberbach1990}; recent work has moved toward legislator behavior within oversight venues. \citet{eldes2023} analyze congressional oversight hearings and distinguish two currents in member questioning: the gathering of information and the staging of confrontation. Their central finding is that confrontation is patterned by partisan alignment, with members of the party opposed to the administration confronting witnesses at higher rates. The finding concerns behavior within a single venue. The natural extension, and the first theoretical position this article tests, is that confrontation crosses venues: a member who confronts an administration figure in one setting acquires commitments that shape behavior in later settings.

Continuity could arise through consistency pressure, through reputational investment, or through the economics of expertise. Consistency pressure operates because a legislator who declared a nominee unfit has issued a public prediction that the ministry will be led badly; subsequent scrutiny of that ministry is the natural way to vindicate the prediction, and inattention would expose the original opposition as costless talk. Reputational investment operates because opposition at a nationally televised hearing is a position-taking act in the classic sense of \citet{mayhew1974}, and position taking generates returns only if audiences can connect the position to later behavior; a legislator known as the nominee's chief antagonist plausibly has a branded issue in that ministry's performance. The expertise mechanism is humbler: preparing to interrogate a nominee requires acquiring knowledge of the ministry's programs and failures, and expertise lowers the cost of future questioning, so once the fixed cost is sunk, the marginal audit question directed at that ministry is cheap. All three mechanisms point the same way. Extended to the appointment context, they predict that a legislator who opposed a confirmed nominee should allocate a durably larger share of oversight attention to the confirmed ministry than the legislator did before the hearing.

\subsection{Confirmation Hearings as Party Theater}

The Korean literature on confirmation hearings supports a different reading of what hearing opposition is. Content analyses of prime ministerial hearings find that questioning tracks the state of executive-legislative relations rather than the properties of individual nominees, with hearings functioning as a stage on which the standing conflict between government and opposition is performed \citep{choi2008}. Topic-model analyses of hearing transcripts reach a compatible conclusion: the discourse of the hearing is structured by the government-assembly relationship of the moment \citep{jeon2018}. Studies of confirmation outcomes find that approval and rejection are driven by the partisan configuration surrounding the nomination rather than by nominee quality as such \citep{yoon2020}. Across these studies the individual legislator appears as the bearer of a party role. Opposition to the nominee is assigned, performed for the hearing day, and discharged when the hearing adjourns.

If hearing opposition is a discharged role, it should leave no residue in the legislator's subsequent behavior. The audit is governed by a different routine: committee jurisdiction determines which ministries a legislator can question, party strategy determines the topics of the season, and the individual position taken months earlier at a hearing determines nothing. The comparative literature on parliamentary questions reinforces this expectation from a different direction \citep{martin2011}. Opposition parties use questions for selective scrutiny tied to the strategic issue opportunities of the moment, not to standing grudges \citep{senninger2017}. Question agendas respond to media attention and focusing events, and the response is concentrated in the immediate window \citep{walgrave2007, vliegenthart2016}. Nothing in this literature suggests that a months-old appointment fight should organize a legislator's question allocation at the next audit.

\subsection{Attention Dynamics and the Audit Calendar}

The timing structure of the Korean case disciplines what either theory can claim. The state audit is not a continuous oversight channel; it is a calendar-fixed annual burst embedded in the National Assembly's rigid institutional time structure, in which fixed session windows govern when legislative attention can be spent at all \citep{ka2025}. Ministerial hearing cohorts concentrate in spring and early summer, so the audit arrives five to ten months after the typical hearing. The political-attention literature speaks directly to what should survive such a gap. Media effects on parliamentary agendas operate at the weekly-to-monthly scale and dissipate quickly \citep{walgrave2007}; pooled evidence from seven countries confirms that media-to-questions transfer is strongest for opposition parties and concentrated in the immediate window \citep{vliegenthart2016}; and even major focusing events produce attention spikes that decay within months as agendas revert to baseline \citep{birkland1998}.

The audit is therefore a demanding but fair venue for the continuity prediction. Demanding, because the decay literature suggests any confrontation-generated attention could be fully discharged before October; fair, because the carry-over claim in the Korean reform debate names the audit specifically, and the argument that hearings improve oversight presumes the improvement shows up where oversight institutionally happens. A further implication concerns the scope of any finding. A null at the audit does not rule out carry-over at shorter range, for instance in current-issues questioning during the weeks after appointment. The claim tested here is the claim the debate makes, and the decay literature independently predicts that the five-to-ten-month gap is long enough for full dissipation of short-lived attention effects.

\subsection{Hypotheses and Pre-Registered Thresholds}

The two theoretical positions yield predictions about a single quantity: the change, from the audit before the hearing to the audit after it, in the share of a legislator's audit questions directed at the confirmed ministry, compared between legislators who opposed and legislators who supported the nominee.

\begin{quote}
\textbf{H1 (position-taking continuity).} Legislators who opposed a confirmed nominee increase the confirmed ministry's share of their audit questions relative to supportive legislators on the same committee. The pre-registered threshold for a substantively meaningful effect is five percentage points.
\end{quote}

\begin{quote}
\textbf{H2 (party theater).} The difference-in-differences lies within the pre-registered equivalence margin of zero. Audit allocation reverts to committee jurisdiction and party role once the hearing ends.
\end{quote}

Two secondary margins were specified in advance, each with its own decision rule fixed before estimation. The intensity margin asks whether continuity, if diluted by pooling all opposed legislators, surfaces among those who engaged the nominee most heavily: within opposition questioners, a one-standard-deviation increase in hearing engagement, measured as the count of question dyads the legislator directed at the nominee and standardized within committee, must raise the difference-in-differences outcome by at least 2.5 percentage points, half the original threshold, for continuity to survive in diluted form. The levels margin arose from early analysis, in which opposed legislators appeared to devote more of their audit questions to the ministry already before the hearing, a pattern initially read as favoring the party-theater account, since it suggests opposition assignment and ministry attention are parallel expressions of the same party role. The pre-committed rule stated that this levels gap would count as a finding only if it survived committee fixed effects; if it collapsed, it would be reported as composition and dropped from the adjudication. Both rules were honored, and the second fired against the project's own earlier reading, as the Results section reports.

Finally, the design commits to the standards of the negligible-effect literature. A null claim is informative only against an explicitly declared magnitude, demonstrated by confidence intervals that exclude effects larger than the declared threshold \citep{rainey2014}, and placebo claims should be stated as equivalence results rather than as failures to reject \citep{hartman2018}. The five-percentage-point threshold, the 2.5-point diluted bar, and the placebo margins reported below were all fixed before the corresponding estimates were computed.

\section{Data and Method}

\subsection{Data}

The analysis draws on the hearings module of the Korean National Assembly (KNA) database, a corpus of committee transcripts spanning the 16th through 22nd Assemblies that contains 9.9 million individual speech acts and 7.4 million question-and-answer dyads linking a questioning legislator to an answering witness. Within this corpus I identify, for each ministerial confirmation hearing, the set of legislators who questioned the nominee, and I follow those same legislators into the annual state audits immediately preceding and immediately following the hearing. The unit of analysis is the legislator-ministry pair: a legislator who questioned a nominee for a given ministry, observed in both the pre-hearing and post-hearing audits of the standing committee holding jurisdiction over that ministry.

Three cohorts of ministerial nominees enter the analysis. Cohort 1 comprises nominees heard in 2020 and 2021, with the audits of October 2020 and October 2021 as the pre- and post-periods. Cohort 2 comprises nominees heard in 2023, bracketed by the audits of 2022 and 2023. Cohorts 1 and 2 together constitute the pooled pre-registered sample of 278 legislator-ministry pairs across 191 legislators. Cohort 3 comprises the nominees of May 2022, whose pre- and post-audits span the transfer of power from the Moon administration to the Yoon administration; because the governing status of every party flips between its two audits, this cohort serves as an out-of-sample consistency check under a full change of government rather than as part of the pooled test. Table~\ref{tab:desc} summarizes the samples.

\begin{table}[!htbp]
\centering
\caption{Analysis Samples and Descriptive Quantities}
\label{tab:desc}
\begin{tabular}{lc}
\toprule
Quantity & Value \\
\midrule
Question-and-answer dyads, full corpus (16th to 22nd Assemblies) & approx. 7.4 million \\
Speech acts, full corpus & approx. 9.9 million \\
Hand-coded hearing roster, cohorts 1 and 2 (members) & 400 \\
Hand-coded hearing roster, cohort 3 (members) & 325 \\
Legislator-ministry pairs, pooled sample (cohorts 1 and 2) & 278 \\
\quad Legislator clusters & 191 \\
\quad Cohort 1 pairs (2020 to 2021 hearings) & 194 \\
\quad Cohort 2 pairs (2023 hearings) & 84 \\
Legislator-ministry pairs, cohort 3 (2022 nominees) & 125 \\
\quad Supportive pairs & 50 \\
\quad Opposed pairs & 75 \\
Within-opposition dose sample, pooled (pairs) & 138 \\
Retention into both audits, opposed pairs (share of 132) & 0.864 \\
Retention into both audits, supportive pairs (share of 197) & 0.832 \\
Hearing questions to nominee, cohort-1 opposition (median) & 31 \\
\bottomrule
\multicolumn{2}{l}{\footnotesize \parbox{0.92\textwidth}{Pairs are legislator-ministry combinations among hearing questioners observed in both the pre- and post-hearing audits. Corpus totals are reported to two significant figures and marked as approximate. The interquartile range of the hearing-dose count in the cohort-1 opposition sample runs from 17 to 51 questions.}} \\
\end{tabular}
\end{table}

\textit{Treatment.} A legislator-ministry pair is coded opposed when the legislator questioned the nominee and belonged, at the hearing date, to a party opposed to the nominating president; pairs of legislators from the president's party are coded supportive. The construct this primary coding measures deserves explicit statement. Party-based opposition status captures the assigned opposition role at the hearing, not an individually expressed position, so it tests the party-theater account on its own terms while testing the individual position taking that H1 invokes only indirectly. The analysis therefore also codes opposition directly from the legislator's own hearing speech, classifying a legislator as opposed when the legislator called for the nominee's withdrawal or declared the nominee unfit; this speech-based coding is the construct-faithful measure of individually expressed opposition, and it is reported alongside the primary specification in the robustness analysis, where it returns the same null. Hearing rosters were hand-coded before any outcome analysis, with 400 member entries for cohorts 1 and 2 and 325 for cohort 3 written to append-only dictionaries. Two coding hazards deserve note because both were caught by internal checks. The first concerns party labels: the corpus records each legislator's party as a term-start snapshot, so members of parties that renamed or merged mid-term carry outdated labels; coding governing status from raw labels inverts 88 of 97 codings in one cohort, and all analyses therefore recode party by the nominating president's party at the speech date, predecessor labels included. The second concerns identity: a name-keyed merge contaminated one supportive-side observation because the 21st Assembly seats two legislators who share a name; all merges in the replication package key on unique member identifiers with a uniqueness assertion, and the corrected merge leaves every headline quantity unchanged at reporting precision.

\textit{Outcome.} The outcome is the confirmed ministry's share of the legislator's question dyads at the state audit, differenced between the post- and pre-hearing audits and expressed in percentage points. The placebo outcome applies the same construction to named government agencies that no nominee in the cohort headed; if opposed legislators simply become more active questioners of government witnesses in general, the placebo moves with the main outcome, whereas a ministry-directed carry-over moves only the main outcome. The intensity measure, used in the within-opposition dose test, is the total count of hearing question dyads the legislator directed at the ministry's nominee or nominees, standardized within committee because hearing length varies mechanically across committees.

\textit{Reproducibility.} The full pipeline reruns end to end deterministically: all persisted artifacts reproduce byte-identically, and the hand-coded dictionaries reproduce the analysis samples exactly on the unique-identifier key. All estimates reported below were regenerated in the final consolidation run rather than transcribed from earlier analyses.

\subsection{Identification Strategy}

The primary specification is a difference-in-differences model estimated in changes. For legislator $i$ questioning ministry $m$ under the jurisdiction of committee $c$,
\begin{equation}
\Delta Y_{im} = \beta_1 \, \text{Opposed}_{im} + \delta_{c(i,m)} + \epsilon_{im},
\label{eq:main}
\end{equation}
where $\Delta Y_{im}$ is the change in ministry $m$'s share of legislator $i$'s audit questions from the pre-hearing to the post-hearing audit, $\text{Opposed}_{im}$ indicates hearing opposition, and $\delta_{c(i,m)}$ are committee fixed effects. Standard errors are clustered by legislator throughout. Because the model is estimated in changes, any time-invariant difference between opposed and supportive legislators in their level of attention to the ministry is removed by construction; the committee fixed effects additionally absorb jurisdiction-specific shifts in audit agendas between the two years. The within-opposition dose test replaces the treatment indicator with the standardized engagement measure,
\begin{equation}
\Delta Y_{im} = \beta_2 \, \text{Dose}_{im} + \delta_{c(i,m)} + \epsilon_{im},
\label{eq:dose}
\end{equation}
estimated on opposition questioners only.

Four threats to inference structure the design. Nonrandom assignment comes first: opposition is not randomly assigned but follows party. Nonrandom assignment, however, is not a nuisance but part of the estimand: the question is whether the opposition role, however assigned, generates individual carry-over, and the changes specification compares opposed and supportive members of the same committee across the same pair of audits. The identifying assumption is parallel trends in ministry shares absent the hearing, and the placebo-agency outcome provides a direct check, since any committee-wide or party-wide trend in questioning activity should appear there as well. Differential attrition is the second threat: selective exit could distort the sample of pairs observed in both audits, but retention is closely balanced, at 0.864 for opposed and 0.832 for supportive pairs. Mechanical composition is the third, because composition can masquerade as behavior: in committees whose audit witness list is dominated by a single ministry, the share outcome is pinned near its ceiling, and one such committee, the Gender Equality and Family Committee, has a mean pre-period share of 0.963; the levels analysis below shows how this cell alone manufactured an apparent attention gap. Power is the fourth, since the informativeness of a null depends on it. The pooled design's minimum detectable effect at 80 percent power is 3.1 percentage points, comfortably below the pre-registered five-point threshold, so the design binds: an effect of the size the continuity hypothesis predicts would have been detected with at least 80 percent probability. Following \citet{rainey2014}, the null claim is stated against the declared threshold, and following \citet{hartman2018}, placebo and main estimates are additionally subjected to two-one-sided equivalence tests at pre-specified margins.

The dose test carries one further caveat that no specification can remove. Legislators who question a nominee intensively may be the committee's pre-existing specialists in that ministry, so the dose partly proxies standing portfolio interest rather than hearing-generated animus. The changes specification absorbs the level of specialization but not differential trends among specialists. For this reason, and because the achievable opposition sample cannot resolve effects of the observed magnitude, the dose result is reported as a bound rather than as a finding.

\section{Results}

\subsection{Main Estimates}

Table~\ref{tab:main} reports the difference-in-differences estimates of Equation~\ref{eq:main}. Column 1 gives the pooled pre-registered test; columns 2 and 3 disaggregate by cohort; column 4 reports the out-of-sample government-change cohort; column 5 replaces the outcome with the placebo-agency share.

\begin{table}[!htbp]
\centering
\caption{Confirmation Hearing Opposition and Audit Question Allocation}
\label{tab:main}
\begin{tabular}{lccccc}
\toprule
& (1) & (2) & (3) & (4) & (5) \\
& Pooled & Cohort 1 & Cohort 2 & Cohort 3 & Placebo \\
\midrule
Opposed & $-0.86$ & $-1.10$ & $1.30$ & $0.09$ & $-0.85$ \\
& (1.11) & (1.28) & (2.42) & (1.99) & (0.94) \\
\midrule
N (legislator-ministry pairs) & 278 & 194 & 84 & 125 & 278 \\
Legislator clusters & 191 & 150 & -- & 87 & 191 \\
Committee FE & Yes & Yes & Yes & Yes & Yes \\
MDE (80\% power, pp) & 3.1 & -- & -- & -- & -- \\
\bottomrule
\multicolumn{6}{l}{\footnotesize $^{*}p<0.10$, $^{**}p<0.05$, $^{***}p<0.01$. SE clustered by legislator in parentheses.} \\
\multicolumn{6}{l}{\footnotesize \parbox{0.92\textwidth}{Outcome: change in the confirmed ministry's share of the legislator's audit questions (percentage points). Column 5 outcome: change in the share directed at named non-confirmed agencies. Cohort 3 spans the 2022 change of government between its two audits and is reported as an out-of-sample consistency check. No estimate is significant at conventional levels. Dashes mark quantities not separately computed for that column: the minimum detectable effect is a design quantity defined for the pooled pre-registered test, and the cohort-2 cluster count is not reported separately.}} \\
\end{tabular}
\end{table}

The pooled estimate is a decline of 0.9 percentage points in the confirmed ministry's share, with a 95 percent confidence interval running from roughly $-3.0$ to $+1.3$ percentage points (Table~\ref{tab:main}, column 1). Nothing in a legislator's audit allocation shifts toward the ministry whose nominee the legislator opposed. The estimate is not merely non-significant; it is signed against the continuity prediction, and the design's minimum detectable effect of 3.1 percentage points binds well below the pre-registered five-point threshold, so the pre-set falsification condition for H1 is met. The cohort-level estimates in columns 2 and 3 are individually consistent with the pooled null, and the government-change cohort in column 4 is null as well: even when the opposed legislators moved from governing party to opposition between their two audits, their allocation toward the contested ministry did not move. The placebo estimate in column 5 is the most instructive single number in the table. Its point estimate differs from the main estimate by less than 0.01 percentage points, indicating that whatever small common drift exists in questioning activity is wholly generic rather than directed at the confirmed ministry.

Figure~\ref{fig:coef} displays the five estimates of Table~\ref{tab:main}, together with the own-speech robustness coding discussed below, against the pre-registered threshold band.

% R code for Figure 1 executed automatically (see figures/fig_1.R)

\begin{figure}[!htbp]
\centering

\includegraphics[width=\textwidth]{figures/2026-08-24_r27/fig_1.pdf}
\caption{Difference-in-differences estimates across samples and specifications. Points are estimates in percentage points; horizontal lines are 95 percent confidence intervals; dotted vertical lines mark the pre-registered five-point threshold.}
\label{fig:coef}
\end{figure}

\subsection{Equivalence Tests}

A null result is informative only if the data affirmatively rule out effects of the declared size \citep{rainey2014}. Figure~\ref{fig:tost} presents the two-one-sided equivalence tests, plotting the 90 percent confidence intervals for the main and placebo estimates against the pre-registered five-point margin and the stricter 2.5-point margin.

% R code for Figure 2 executed automatically (see figures/fig_2.R)

\begin{figure}[!htbp]
\centering
\includegraphics[width=\textwidth]{figures/2026-08-24_r27/fig_2.pdf}
\caption{Two-one-sided equivalence tests at the pre-registered and half-registered margins. Intervals are 90 percent confidence intervals, following the equivalence-testing convention.}
\label{fig:tost}
\end{figure}

The results are asymmetric by design, and the asymmetry should be stated exactly. The placebo estimate is formally equivalent to zero within plus or minus 2.5 percentage points: its 90 percent interval sits entirely inside even the half-registered margin, which converts the informal claim that placebo agencies move identically into the affirmative equivalence statement that \citet{hartman2018} recommend. The main estimate is formally equivalent to zero within the pre-registered five-point margin, but not within the stricter 2.5-point margin, because its lower bound extends slightly past that line. The correct summary is therefore that the effect of hearing opposition is bounded below the substantive threshold declared before estimation, while the placebo is equivalent to zero at even half that threshold. Claiming the stricter equivalence for the main effect would overstate the evidence, and the paper does not.

\subsection{Robustness}

Table~\ref{tab:robust} probes the null against the two remaining researcher degrees of freedom: the coding of treatment and the construction of the outcome.

\begin{table}[!htbp]
\centering
\caption{Robustness: Alternative Treatment Coding and Alternative Outcomes}
\label{tab:robust}
\begin{tabular}{lccc}
\toprule
& (1) & (2) & (3) \\
& Own-speech coding & Log question count & Named-agency share \\
\midrule
Opposed & $-1.60$ & $-0.041$ & $0.20$ \\
& (1.66) & (0.080) & (2.02) \\
\midrule
N (legislator-ministry pairs) & 278 & 278 & 278 \\
Legislator clusters & 191 & 191 & 191 \\
Committee FE & Yes & Yes & Yes \\
\bottomrule
\multicolumn{4}{l}{\footnotesize $^{*}p<0.10$, $^{**}p<0.05$, $^{***}p<0.01$. SE clustered by legislator in parentheses.} \\
\multicolumn{4}{l}{\footnotesize \parbox{0.9\textwidth}{Column 1 codes opposition from the legislator's own hearing speech (calls for withdrawal or disqualification) rather than from party; outcome as in Table~\ref{tab:main}, percentage points. Column 2 outcome: change in the log count of questions to the confirmed ministry (log points). Column 3 outcome: change in the confirmed ministry's share computed among named-agency questions only (percentage points). No estimate is significant at conventional levels.}} \\
\end{tabular}
\end{table}

The null survives every variation. Coding opposition from what legislators actually said at the hearing, the construct-faithful measure of individually expressed position taking, rather than from the party they belonged to, leaves the estimate negative and indistinguishable from zero (Table~\ref{tab:robust}, column 1). Replacing the share outcome with the log count of questions directed at the ministry, which guards against the possibility that shares conceal offsetting movements in overall questioning volume, produces the same conclusion (column 2), as does restricting the denominator to questions aimed at named agencies (column 3). No specification in Table~\ref{tab:main} or Table~\ref{tab:robust} produces a positive estimate that approaches conventional significance, and none approaches the pre-registered threshold.

\subsection{The Levels Margin: An Apparent Gap Dissolves}

An earlier stage of this project identified what looked like the arc's one positive finding: in the raw cohort-1 data, opposed legislators already devoted 4.6 percentage points more of their audit questions to the ministry than supportive legislators did before the hearing ever took place. That pattern was initially read as evidence for the party-theater account, on the logic that opposition assignment and ministry-directed attention are parallel expressions of the same party role. The pre-committed decision rule stated that the gap would count as a finding only if it survived committee fixed effects. Table~\ref{tab:levels} reports the test.

\begin{table}[!htbp]
\centering
\caption{Baseline Levels Gap Under Committee Fixed Effects}
\label{tab:levels}
\begin{tabular}{lccc}
\toprule
& (1) & (2) & (3) \\
& Cohort 1 & Pooled & Pooled, placebo agencies \\
\midrule
Opposed & $0.69$ & $1.53^{*}$ & $2.23^{**}$ \\
& (0.86) & (0.81) & (1.05) \\
\midrule
N (legislator-ministry pairs) & 194 & 278 & 278 \\
Committee FE & Yes & Yes & Yes \\
\bottomrule
\multicolumn{4}{l}{\footnotesize $^{*}p<0.10$, $^{**}p<0.05$, $^{***}p<0.01$. SE clustered by legislator in parentheses.} \\
\multicolumn{4}{l}{\footnotesize \parbox{0.9\textwidth}{Outcome: pre-hearing share of the legislator's audit questions directed at the confirmed ministry (columns 1 and 2) or at named non-confirmed agencies (column 3), in percentage points. The raw cohort-1 gap without fixed effects is 4.61 points; excluding the single-ministry Gender Equality and Family Committee (mean pre-period share 0.963, seven of 11 pairs opposed) it falls to 0.61 points.}} \\
\end{tabular}
\end{table}

The gap does not survive. Holding committee constant, the cohort-1 estimate collapses from 4.6 percentage points to well under one point, with a confidence interval straddling zero (Table~\ref{tab:levels}, column 1). The decomposition identifies the source precisely: one committee does nearly all the work. In the Gender Equality and Family Committee, the confirmed ministry is effectively the entire audit witness list, so every member's share sits near its ceiling, and seven of the committee's 11 pairs happen to be opposed; dropping that single cell reduces the raw gap by roughly seven-eighths. The pooled residual in column 2 is marginal, and column 3 shows why it cannot be read as ministry-directed attention: opposed legislators exhibit a same-signed and larger gap on placebo agencies that no nominee headed. What little remains is the generic pattern that opposition members direct more questions at named government witnesses of every kind. The pre-committed rule therefore fired: the levels register drops from the adjudication, and the simplest description the data license is that the level of audit attention follows committee jurisdiction.

\subsection{The Intensity Margin: A Bounded Residual}

The final margin on which the continuity hypothesis could survive is intensity. If carry-over exists but is diluted by pooling every opposed legislator, it should concentrate among those who engaged the nominee most heavily. Table~\ref{tab:dose} reports the within-opposition dose specifications of Equation~\ref{eq:dose}, and Figure~\ref{fig:dose} displays them against the pre-set diluted-continuity bar of 2.5 percentage points per standard deviation.

\begin{table}[!htbp]
\centering
\caption{Within-Opposition Dose Test: Hearing Engagement and Audit Reallocation}
\label{tab:dose}
\begin{tabular}{lccccc}
\toprule
& (1) & (2) & (3) & (4) & (5) \\
& Pooled & Cohort 1 & Log dose & Raw count & Placebo \\
\midrule
Hearing dose & $1.28^{*}$ & $1.78^{*}$ & $1.34^{**}$ & $0.11^{**}$ & $0.21$ \\
& (0.72) & (0.97) & (0.65) & (0.06) & (0.64) \\
\midrule
N (opposition pairs) & 138 & 80 & 138 & 80 & 138 \\
Legislator clusters & 113 & 60 & 113 & 60 & 113 \\
Committee FE & Yes & Yes & Yes & Yes & Yes \\
\bottomrule
\multicolumn{6}{l}{\footnotesize $^{*}p<0.10$, $^{**}p<0.05$, $^{***}p<0.01$. SE clustered by legislator in parentheses.} \\
\multicolumn{6}{l}{\footnotesize \parbox{0.94\textwidth}{Outcome: change in the confirmed ministry's share of audit questions (percentage points). Dose is the count of hearing question dyads directed at the nominee, standardized within committee in columns 1, 2, and 5; column 3 uses the log of the count; column 4 uses the raw per-question count in cohort 1 (the pooled raw-count slope does not approach significance). Column 5 outcome: change in the placebo-agency share. Excluding the single-ministry committee, the pooled standardized slope is 1.10 (SE 0.76); a top-tercile engagement indicator returns a null.}} \\
\end{tabular}
\end{table}

% R code for Figure 3 executed automatically (see figures/fig_3.R)

\begin{figure}[!htbp]
\centering
\includegraphics[width=\textwidth]{figures/2026-08-24_r27/fig_3.pdf}
\caption{Within-opposition dose slopes against the pre-set diluted-continuity bar. Points are slopes in percentage points per within-committee standard deviation of hearing engagement (per log unit or per question in the alternative measures); the dotted line marks the 2.5-point bar.}
\label{fig:dose}
\end{figure}

The dose results land between the two theories, and the honest classification is a bound. Every specification returns a positive slope, and the placebo slope in column 5 is centered on zero, so whatever tilt exists does not appear on agencies the hearings did not touch. Yet no point estimate reaches the pre-set diluted-continuity bar, and only one of six specifications, the log-dose variant, excludes zero, a pattern consistent with fragility rather than signal. Neither pre-registered decision rule can fire: the survival bar of 2.5 points per standard deviation is not met by any point estimate, while the failure condition, an interval excluding that bar on a sample of at least 150 opposition pairs, cannot fire at the achievable 138. The design's minimum detectable effect on this sample is roughly two percentage points per standard deviation, so an effect of the observed magnitude is unresolvable by construction. What the data support is the bounding statement: at 95 percent confidence, the within-opposition intensity slope lies below 2.7 percentage points per standard deviation of hearing engagement (Table~\ref{tab:dose}, column 1), below half the original threshold. The continuity hypothesis is weakened at this margin, not killed, and the specialization confound discussed above gives an additional reason not to promote the residual tilt to a finding.

\subsection{Summary of Results}

Across 18 pre-committed tests and robustness checks, the pattern is uniform. The pre-registered carry-over prediction fails with a binding minimum detectable effect and a placebo that moves identically; the null is formally equivalent to zero at the registered margin; it holds under an alternative treatment coding, alternative outcomes, and a full change of government; the one apparent positive finding dissolves into committee composition under a pre-committed rule; and the intensity margin is inconclusive rather than resolved, bounded below half the declared threshold. Audit attention, in levels and in changes, tracks committee jurisdiction.

\section{Discussion}

Two literatures made opposite predictions for the same quantity, and the data adjudicate cleanly. Position-taking continuity, extended from \citet{mayhew1974} through the confrontation-in-oversight work of \citet{eldes2023}, predicted that a legislator who staked a public position against a nominee acquires a durable stake in that ministry's failure, raising the ministry's share of the legislator's questions at the following state audit. The Korean confirmation-hearing literature predicted instead that hearing opposition is a party-assigned role for the hearing day, discharged when the hearing ends, with audit allocation reverting to committee jurisdiction \citep{choi2008, jeon2018, yoon2020}. The continuity prediction fails at the party level, with an estimate equivalent to zero at the threshold declared before estimation, and it fails again under a full change of government between audits. Within opposition questioners, the intensity margin is bounded below half the pre-set bar. What the levels of audit attention track is neither hearing opposition nor its intensity but committee jurisdiction: the apparent baseline gap between opposed and supportive legislators dissolves under committee fixed effects and is matched by a larger placebo gap. Confirmation conflict and audit allocation appear to be two separate institutional routines, not a causal chain.

The adjudication is narrower than a simple victory for the party-theater account, and the narrowing deserves emphasis. The zero-change prediction of the party-theater reading held, but the reading had also seemed to gain support from the baseline levels pattern, and that support is now gone: the levels evidence turned out to discriminate for nothing except the composition of committee witness lists. The description the data license is more austere than either theory. Allocation follows jurisdiction. Party assignment may well govern who performs opposition at the hearing, as the Korean literature documents, but the performance leaves no measurable residue in the performer's subsequent allocation of oversight attention, and the party role does not even generate the standing attention premium that a strong role-based account might expect.

The results align well with what the political-attention literature would predict for this venue. Attention transfer into parliamentary questions is fast and short-lived, operating at the weekly-to-monthly scale \citep{walgrave2007, vliegenthart2016}, and even major focusing events produce spikes that revert within months \citep{birkland1998}. The state audit sits five to ten months downstream of the typical hearing, inside a calendar that fixes when oversight attention can be spent at all \citep{ka2025}. A carry-over effect could in principle exist at shorter range, in current-issues questioning during the weeks after appointment, and be fully discharged by October. This study does not test that channel, and the null here does not speak to it. What the study shows is that confirmation conflict does not carry into the audit, the oversight venue where the Korean reform debate has assumed the hearing's externality lands. The scope statement and the venue were both fixed by the prior being tested, not chosen after the fact.

The results also refine, rather than contradict, the confrontation findings that motivated the continuity hypothesis. \citet{eldes2023} document that partisanship shapes the style of oversight questioning, with opposition-aligned members confronting witnesses at higher rates. Nothing here disputes that. The government-change cohort is instructive on exactly this point: when the opposed legislators' party moved from governing to opposition between the two audits, a shift that should transform confrontation incentives, their allocation of questions across ministries still did not move. Partisan alignment appears to govern how legislators question, while jurisdiction governs whom they question; the appointment fight changes neither. The finding similarly complements \citet{senninger2017}: opposition scrutiny follows the strategic opportunities of the moment, and a months-old hearing position does not seem to constitute one.

Several limitations bound these conclusions, and they should be stated plainly. The outcome is allocation, not content: legislators who opposed a nominee might question the confirmed ministry no more often but more aggressively, and a tone-based or confrontation-coded outcome could in principle detect a carry-over that share-based allocation misses; that channel remains untested here. The intensity margin is unresolved rather than resolved: the dose slope is positive in every specification and clean on placebo agencies, but the achievable opposition sample cannot distinguish an effect of the observed size from zero, and the specialization confound means even a resolved positive slope would not identify hearing-generated attention, so the bound reported is the honest object. Treatment is observational, and the primary coding is party-based: opposition follows party, the party-based indicator captures the assigned role rather than an individually expressed position, and although the changes specification, committee fixed effects, placebo outcomes, balanced retention, and the speech-based robustness coding address the sharpest threats, no design element delivers the counterfactual of a randomly assigned hearing antagonist. Some cells are thin: the cohort-2 supportive side contributes 26 pairs, and single-ministry committees are unusable for share outcomes, a design lesson this study learned by initially mistaking one such committee's composition for behavior. The setting, finally, is one country with a distinctive oversight calendar: the tight coupling of hearings and audits to the same standing committees is what makes the test possible, but it also means the decay window between conflict and oversight opportunity is institutionally fixed at five to ten months, and systems with continuous questioning channels might show short-range dynamics this design cannot observe.

Finally, the trajectory of the project is part of the evidence, and it is reported as procedure rather than confession. Two findings were overturned by pre-committed decision rules during the research: the original carry-over prior fell to its pre-registered falsifier, and the baseline levels gap, at one stage the project's candidate positive finding, fell to the fixed-effects rule that had been attached to it in advance. Both retreats were logged when they occurred, and neither headline claim depends on any result that failed its own pre-set test. In the framework of \citet{rainey2014}, the value of a null rests on declared magnitudes and demonstrated power; in this study it rests additionally on decision rules that fired against the analysts' own earlier readings, which may be the more demanding certification.

\section{Conclusion}

This article asked whether the conflict legislators stage over ministerial appointments carries into their subsequent oversight of the ministries those appointments fill. Adjudicating between a position-taking continuity extension of the oversight literature and the party-theater reading of Korean confirmation hearings, I linked hand-coded hearing rosters and party-based opposition status for three nominee cohorts, supplemented by a speech-based coding of individually expressed opposition, to 7.4 million question-and-answer dyads from the Korean National Assembly and estimated a pre-registered difference-in-differences model of audit question allocation. The answer is a null with unusually complete credentials: equivalent to zero at the threshold declared before estimation, placebo-verified in the formal equivalence sense, robust to alternative treatment codings and outcomes, stable across a change of government, and accompanied by a bounded intensity residual and a dissolved levels gap. Audit attention follows committee jurisdiction; the appointment fight leaves no measurable trace on it.

The contribution is threefold. Theoretically, the finding sets a scope condition on position-taking continuity: whatever stakes a confirmation fight creates do not survive the five-to-ten-month passage to the next audit, consistent with the fast decay the attention literature documents. For the study of the Korean National Assembly, the finding replaces an assumed linkage between two prominent oversight institutions with a measured absence of one. For institutional reform, the implication is direct: proposals to restructure confirmation hearings should not be premised on an oversight externality at the audit, because the carry-over such proposals presume is equivalent to zero at the threshold declared before estimation, and the residual intensity slope, which the design cannot fully resolve, is bounded below half that threshold. Confirmation hearings may serve screening, accountability, or position-taking functions on their own terms; improving later audit scrutiny does not appear to be among them.

Future research should pursue the two channels this design could not reach. The first is short-range carry-over, in current-issues questioning during the weeks immediately after appointment, where the attention literature predicts any effect would concentrate. The second is tone: whether opposed legislators question the confirmed ministry differently, if not more. Comparative extension to systems with continuous questioning channels would test whether the null reported here is a property of appointment conflict generally or of calendars that separate conflict from oversight opportunity by months. Each extension would benefit from the same discipline this study relied on: thresholds declared first, and decision rules that are allowed to fire against the researcher.

\bigskip
\noindent\textit{This working paper was generated by AI research agents as an
experimental output. It has not been peer-reviewed or fact-checked.
Do not cite or use in any academic, policy, or professional context.}

\bibliographystyle{apsr}
\bibliography{references_r27}


\end{document}
